\section{Exact external-response identities}\label{sec:external-response}

We prove the exact finite-volume identities of
\cite[\S20]{QED3Master}. We express the generating function of a curvature observable as a
determinant in an external connection. The identity is first proved
in each integer flux sector. In the zero-flux sector an unfolding
of the integer cochains replaces the compact photon integral by a
Gaussian integral over transverse cochains and an average over the
three twists. We compute its normalization before applying the
source shift.

\subsection{The determinant at complex connections}
\label{subsec:external-complex-determinant}

Use lattice units and fix $q>0$, $r_q=2\pi/q$ and a real lattice
mass $\mu_{\mathrm{lat}}$. Put
$M_{\mathrm{ch}}=\mu_{\mathrm{lat}}+\frac{\rho}{2}\sum_iw_i$.
For the law in \cref{thm:rrm-main},
$\mu_{\mathrm{lat}}=\boldsymbol\mu_N$.
The weights are the standard chamber weights, and $0<\rho<1$.
For an integer triple $n=(n_1,n_2,n_3)$ let
\[
 V=|n|=n_1n_2n_3,\qquad
 B_n=\prod_{\mu=1}^3\{0,\ldots,n_\mu-1\},\qquad
 \mathcal L_n=B_n\times\{1,2,3\}.
\]
In this subsection each $n_\mu$ may be as small as one.
Cochains are periodic functions on the oriented covering cells,
so the $3V$ link orbits remain distinct even at those small sides.
The chamber spinors are antiperiodic in each direction and have
dimension $192V$. Write $D_B$ for the massless chamber operator
with phases $\exp(iqB)$, and $\det_n$ for its finite-dimensional
determinant. In the phase convention of
\cref{subsec:rrm-regulator}, $D_B=D^{\mathrm{ch}}[qB]$.
The weight $1/48$ in the spinor pairing does not rescale this
determinant.

The determinant is a polynomial in the $6V$ variables
$\exp(iqB(b))$ and $\exp(-iqB(b))$, of degree at most $32$
in each. In particular it is entire in the complex link values.
To see the degree bound, choose the site basis on
$B_n\times O_{\mathrm{cub}}$ and include the antiperiodic sign
whenever a neighbor is reduced into $B_n$. Each of the four
rows at a site $(c,g)$ has constant entries except for the
one external-hop block. If $ge_1=e_\mu$, this block contains
$\exp(iqB_\mu(c))$; if $ge_1=-e_\mu$, it contains
$\exp(-iqB_\mu(c-e_\mu))$. For a fixed oriented photon link
there are eight labels at each end. Thus each phase variable
occurs in exactly $8\cdot4=32$ rows. The Leibniz determinant
chooses one entry in each row, proving the bound. The five
site labels of the diagonal and the four hops are distinct,
even if a side is one, so no other phase is hidden in one
of these blocks.

The two invariances needed below hold for complex cochains:
\begin{equation}\label{eq:external-determinant-invariances}
 \det_n(D_{B+r_qk}+\mu_{\mathrm{lat}})
       =\det_n(D_B+\mu_{\mathrm{lat}}),\qquad
 \det_n(D_{B+d\lambda}+\mu_{\mathrm{lat}})
       =\det_n(D_B+\mu_{\mathrm{lat}}),
\end{equation}
where $k$ is integer valued and $\lambda$ is a periodic complex
zero-cochain. The first identity follows directly from the phases.
For the second, multiplication by $\exp(iq\lambda(c))$ preserves
antiperiodicity and has inverse multiplication by
$\exp(-iq\lambda(c))$. On an external hop the factor
\[
 \mathrm e^{iq(B-d\lambda)(c,c')}\mathrm e^{iq\lambda(c')}
                                  =\mathrm e^{iq\lambda(c)}\mathrm e^{iqB(c,c')}
\]
gives $D_{B-d\lambda}=\mathrm e^{iq\lambda}D_B\mathrm e^{-iq\lambda}$.
Internal hops lie over one photon site and give the same
conjugation. This proves the determinant identity.

For $s\ge0$ define
\begin{align}
 K_n(s)&=
 \left(|M_{\mathrm{ch}}|+\frac{1+\rho}{2}
           \left(\sum_iw_i-6+6\mathrm e^{qs}\right)\right)^{192V},
           \notag\\
 |\det_n(D_B+\mu_{\mathrm{lat}})|
                   &\le K_n(\|\operatorname{Im}B\|_\infty).
                   \label{eq:external-determinant-strip}
\end{align}
Indeed each chamber bond involution permutes the site representatives.
For an internal hop its phase has modulus one; for an external
hop the modulus is at most $\mathrm e^{q\|\operatorname{Im}B\|_\infty}$.
The hopping matrix has norm $w_i(1+\rho)/2$. The operator norm
is therefore bounded by the base of the power in
\eqref{eq:external-determinant-strip}, and a determinant in dimension
$192V$ is bounded by that norm to the same power.
This estimate is uniform in the real part of the connection.

The twist connection is
\begin{equation}\label{eq:external-twist}
 H_\theta=\sum_\mu\frac{\theta_\mu}{qn_\mu}\mathbf1_\mu,
 \qquad
 \mathbf1_\mu=\sum_{c\in B_n}\mathbf1_{(c,\mu)}.
\end{equation}
Its determinant has period $2\pi$ in each complex twist variable.
For later use we need the cochain implementing this period. Set
\[
 k^{(\mu)}_\nu(c)=\delta_{\mu\nu}1_{\{c_\mu=n_\mu-1\}},
 \qquad
 \lambda^{(\mu)}(c)=-\frac{r_q}{n_\mu}c_\mu,\qquad c\in B_n,
\]
with periodic extension. On links not crossing the $\mu$ seam,
$d\lambda^{(\mu)}=-r_q/n_\mu$ in direction $\mu$ and zero
in the other directions. On a seam link its value is
$r_q-r_q/n_\mu$. Thus
\begin{equation}\label{eq:external-twist-period}
 H_{2\pi e_\mu}=r_qk^{(\mu)}-d\lambda^{(\mu)}.
\end{equation}
The two invariances in
\eqref{eq:external-determinant-invariances} prove the asserted
twist period. For $n_\mu=1$, $\lambda^{(\mu)}=0$ and
$k^{(\mu)}=\mathbf1_\mu$, so the same formula holds.

\subsection{Real and integer cochains on the torus}
\label{subsec:external-cochains}

From now on $n_\mu\ge4$. The counting pairings on cochains
are extended bilinearly when their arguments are complex.
Use
\[
 V_{\mathrm g}=dC^0(n),\qquad
 V_{\mathrm h}=\operatorname{span}_{\R}\{\mathbf1_1,\mathbf1_2,\mathbf1_3\},
 \qquad
 V_{\mathrm t}=(V_{\mathrm g}+V_{\mathrm h})^\perp.
\]
We recall the cochain arguments needed for the unfolding and
for its Jacobian. In particular the integer condition in the
flux sum requires an integer primitive, not merely a real one.

Let
\[
 \mathcal P_n^{\mathrm{per}}=
   \left\{p:p_\mu=2\pi j_\mu/n_\mu,\ 0\le j_\mu<n_\mu\right\},
 \qquad e_p(c)=V^{-1/2}\mathrm e^{ip\cdot c},\qquad
 \xi_\mu(p)=\mathrm e^{ip_\mu}-1.
\]
The characters $e_p$ are orthonormal. At each momentum $p$ the
coboundary on exterior-valued cochains is exterior multiplication
by $\xi(p)$. If $\iota_{\bar\xi}$ denotes contraction by $\bar\xi$,
then
\[
 \iota_{\bar\xi}(\xi\wedge u)+
                  \xi\wedge\iota_{\bar\xi}u=|\xi|^2u.
\]
For $p\ne0$ a closed coefficient $u$ is consequently
$\xi\wedge(|\xi|^{-2}\iota_{\bar\xi}u)$, and is exact.
At zero momentum the cochains are constant and the
coboundary vanishes. Nonzero and zero modes are orthogonal,
and taking real parts of a complex primitive gives a real
primitive for a real cochain. Therefore
\begin{equation}\label{eq:external-Hodge}
 \ker(d:C^1\to C^2)=V_{\mathrm g}\oplus V_{\mathrm h},
 \qquad C^1=V_{\mathrm g}\oplus V_{\mathrm h}\oplus V_{\mathrm t},
 \qquad
 \ker(d:C^2\to C^3)=dC^1\oplus\operatorname{Harm}^2.
\end{equation}
Here $\operatorname{Harm}^2$ is the space of constant two-cochains.
The same Fourier formula gives $d^2=0$.

For a closed integer two-cochain $z$, define
$w_{\mu\nu}(z)$ as its sum over one coordinate $(\mu,\nu)$
two-torus. The sum is independent of the transverse coordinate:
sum $dz=0$ over the layer of cubes between two consecutive
planes. Differences in the two tangential directions cancel
by periodicity; the remaining difference is that of the two
plane sums. The fluxes give
\begin{equation}\label{eq:external-integer-two-exactness}
 z\in dC^1(n;\Z)\quad\Longleftrightarrow\quad
                         w_{\mu\nu}(z)=0\ \text{for all }\mu<\nu.
\end{equation}
The integer coefficients follow from the
cellular cohomology of the product of the three subdivided
circles. Its cochain complex is precisely the integer complex
used here. The three coordinate two-tori give a basis of
$H_2(T^3;\Z)$, and $H_1(T^3;\Z)$ is free. The universal
coefficient theorem identifies $H^2(T^3;\Z)$ with
$\operatorname{Hom}(H_2(T^3;\Z),\Z)$ by evaluation.
The class of $z$ evaluates to $w_{\mu\nu}(z)$ on each basis
two-torus, so it vanishes exactly when these sums vanish.
Vanishing of this cellular class is membership in
$dC^1(n;\Z)$.
\cite[Theorems~3.2, 3.5 and~3B.6]{Hatcher2002}

We next describe the closed integer one-cochains. If $d\lambda$
is integer valued, its sum along a path from $0$ to $c$ shows
that $\lambda(c)-\lambda(0)$ is integer valued. If $d\lambda=0$,
the same path argument shows that $\lambda$ is constant.
It follows that
\[
 C^1(n;\Z)\cap V_{\mathrm g}=dC^0(n;\Z).
\]
Each seam cochain $k^{(\mu)}$ in
\eqref{eq:external-twist-period} is closed, because its
multiple by $r_q$ lies in $V_{\mathrm g}+V_{\mathrm h}$.
For an arbitrary closed integer one-cochain $k$,
\eqref{eq:external-Hodge} gives
$k=d\lambda+\sum_\mu t_\mu\mathbf1_\mu$.
Its sum along a periodic $\mu$ loop is $n_\mu t_\mu$,
an integer $j_\mu$. Subtracting
$\sum_\mu j_\mu k^{(\mu)}$ leaves an integer exact cochain,
so
\begin{equation}\label{eq:external-closed-integer-one}
 k=d\lambda'+\sum_\mu j_\mu k^{(\mu)},\qquad
 \lambda'\in C^0(n;\Z),\quad j\in\Z^3.
\end{equation}
The harmonic projection of $k$ determines $j$, and then
determines $d\lambda'$. Moreover
\begin{equation}\label{eq:external-integer-Hodge}
 r_qk=
   \left(r_qd\lambda'+\sum_\mu j_\mu d\lambda^{(\mu)}\right)
                         +H_{2\pi j}.
\end{equation}
The first term is exact and the second is harmonic. Their
coupling through $j$ will matter when choosing a fundamental
domain.

The scalar Laplacian $\Delta_0=d^*d$ has eigenvalues
$|\xi(p)|^2=\sum_\mu(2-2\cos p_\mu)$ and a one-dimensional
kernel. Write
\begin{equation}\label{eq:external-scalar-determinant}
 \det{}'\Delta_0=\prod_{p\in\mathcal P_n^{\mathrm{per}}\setminus\{0\}}
                       |\xi(p)|^2,\qquad
 \mathcal L(n)=\log\det{}'\Delta_0.
\end{equation}
At each nonzero momentum the transverse one-cochains
are the two-dimensional space $\xi(p)^\perp$.
For $v\in\C^3$,
\[
 \sum_{\mu<\nu}|\xi_\mu v_\nu-\xi_\nu v_\mu|^2
           =|\xi|^2|v|^2-|\bar\xi\cdot v|^2.
\]
Thus $d^*d$ acts on that transverse space as multiplication
by $|\xi(p)|^2$, and
\begin{equation}\label{eq:external-transverse-determinant}
 \dim V_{\mathrm t}=2V-2,\qquad
 \det(d^*d|_{V_{\mathrm t}})=(\det{}'\Delta_0)^2.
\end{equation}
Diagonalizing this positive real form gives its Gaussian integral
\begin{equation}\label{eq:external-transverse-Gaussian}
 \mathcal G_n:=\int_{V_{\mathrm t}}\mathrm e^{-\|dA\|^2/2}\,\dd\sigma(A)
                    =\frac{(2\pi)^{V-1}}{\det{}'\Delta_0}.
\end{equation}
The measure $\sigma$ is Euclidean Lebesgue measure for the
counting pairing, not a coordinate measure chosen along a
nonorthogonal basis.

There is one more determinant needed for the exact directions.
The cochains $d\mathbf1_c$, $c\in B_n\setminus\{0\}$,
form both a real basis of $V_{\mathrm g}$ and an integer
basis of $dC^0(n;\Z)$. Indeed one subtracts the value at
zero from a zero-cochain; independence follows because
the resulting zero-cochain has value zero there and
cannot be a nonzero constant. Their Gram determinant is
\begin{equation}\label{eq:external-exact-Gram}
 \det\big(\langle d\mathbf1_c,d\mathbf1_{c'}\rangle\big)_{c,c'\ne0}
                           =\frac{\det{}'\Delta_0}{V}.
\end{equation}
To prove this formula, let $S$ be the scalar Laplacian matrix
in the site basis. The coefficient of $t$ in $\det(t+S)$
is $\det{}'\Delta_0$, by its eigenvalues. Multilinearity
of the determinant in columns identifies that coefficient
with the sum of the $V$ principal minors obtained by
deleting one row and column. All these minors have the same
determinant, since site translations conjugate $S$ by
permutations. The minor at zero is the Gram matrix above.

\subsection{Unfolding the zero-flux sector}
\label{subsec:external-zero-flux}

Let $\Phi$ be a Borel function on $C^2(n)$ with
$|\Phi(u)|\le C_\Phi \mathrm e^{C_\Phi\|u\|}$.
Let $W$ be a bounded continuous complex function on
$C^1(n)$, invariant under addition of real exact cochains
and of $r_q$ times integer one-cochains. Then
\begin{align}
 &\sum_{\substack{z\in C^2(n;\Z),\,dz=0\\w(z)=0}}
       \int_{[0,r_q)^{\mathcal L_n}}
            \Phi(dA+r_qz)\mathrm e^{-\|dA+r_qz\|^2/2}W(A)\,\dd A
       \notag\\
 &\quad=\mathcal N_n
       \int_{[-\pi,\pi]^3}\frac{\dd\theta}{(2\pi)^3}
       \int_{V_{\mathrm t}}
              \Phi(dA)\mathrm e^{-\|dA\|^2/2}W(A+H_\theta)\,\dd\sigma(A),
                                     \label{eq:external-zero-unfolding}\\
 &\mathcal N_n=r_q^{V+2}(\det{}'\Delta_0)^{1/2}.
                                     \label{eq:external-zero-normalization}
\end{align}
Both sides converge absolutely. The coefficient comes from
the lattice of exact gauge transformations and the three
periods of the harmonic directions; we give the unfolding
with both contributions.

First, set
\[
 C_q(\epsilon)=\sup_{y\in\R}\sum_{j\in\Z}
                  \mathrm e^{-\epsilon(y+r_qj)^2}
       \le2+r_q^{-1}\sqrt{\pi/\epsilon}\qquad(\epsilon>0).
\]
The bound follows by keeping the two nearest terms and comparing
the remaining terms with Gaussian integrals over adjacent intervals,
as in \cite[(5)]{ChamberLattice}. Summing each integer coordinate gives
\[
 \sum_{z\in C^2(n;\Z)}\int_{[0,r_q)^{\mathcal L_n}}
              \mathrm e^{-\epsilon\|dA+r_qz\|^2}\,\dd A
                     \le(r_qC_q(\epsilon))^{3V}.
\]
Here the integer closedness and flux restrictions have
been dropped. Completing the scalar square bounds
$|\Phi(u)|\mathrm e^{-\|u\|^2/2}$ by
$C_\Phi \mathrm e^{C_\Phi^2}\mathrm e^{-\|u\|^2/4}$.
Together with boundedness of $W$, this proves absolute
convergence of the left side.

Let $\mathcal Z^1_{\Z}=\ker d\cap C^1(n;\Z)$ and
choose representatives $\mathcal R$ of
$C^1(n;\Z)/\mathcal Z^1_{\Z}$. By
\eqref{eq:external-integer-two-exactness}, $k\mapsto dk$
is a bijection from $\mathcal R$ onto the closed
zero-flux integer two-cochains. Put
\[
 P(A)=\Phi(dA)\mathrm e^{-\|dA\|^2/2}W(A),\qquad
 \mathcal D_1=\bigcup_{k\in\mathcal R}
                     ([0,r_q)^{\mathcal L_n}+r_qk).
\]
The two invariances of $W$ turn the left side of
\eqref{eq:external-zero-unfolding} into
$\int_{\mathcal D_1}P(A)\,\dd A$.
The boxes are disjoint. Every $A$ has a unique decomposition
into its coordinatewise remainder in $[0,r_q)$ and
$r_q$ times an integer cochain. Decomposing that cochain
into its chosen representative and a closed one shows
that $\mathcal D_1$ is a Borel fundamental domain for
$r_q\mathcal Z^1_{\Z}$. The function $P$ is invariant
under that group.

We replace this domain by
\[
 \mathcal F_{\mathrm g}
 =\left\{r_q\sum_{c\ne0}t_c\,d\mathbf1_c:0\le t_c<1\right\},
 \qquad
 \mathcal F_{\mathrm h}=\{H_\theta:\theta\in[-\pi,\pi)^3\},
 \qquad
 \mathcal D_2=\mathcal F_{\mathrm g}+\mathcal F_{\mathrm h}
                                      +V_{\mathrm t}.
\]
To check that it is a fundamental domain for the same
group, write $A=A_{\mathrm g}+A_{\mathrm h}+A_{\mathrm t}$.
Formula \eqref{eq:external-integer-Hodge} shows that
subtracting $r_qk$ changes the harmonic part by
$H_{2\pi j}$ and the exact part by
$\sum_\mu j_\mu d\lambda^{(\mu)}+r_qd\lambda'$.
There is exactly one $j$ placing the harmonic part
in $\mathcal F_{\mathrm h}$. For this $j$ there is
exactly one $d\lambda'$ placing the exact part
in $\mathcal F_{\mathrm g}$, by the integer basis
used in \eqref{eq:external-exact-Gram}. The transverse
part is unchanged. Uniqueness in
\eqref{eq:external-closed-integer-one} gives uniqueness
of $k$. The orthogonal decomposition identifies
$\mathcal D_2$ with a product of Borel sets, so it
is Borel.

For an integrable periodic Borel function, the integral
over two such fundamental domains is the same.
Indeed partition $\mathcal D_1$ into its intersections
with $\mathcal D_2+h$, $h\in r_q\mathcal Z^1_{\Z}$,
and translate each intersection by $-h$. They partition
$\mathcal D_2$, and invariance of both the measure and
the function gives equality. Tonelli proves it first
for nonnegative functions; applying it to the modulus
then justifies the complex case by absolute integrability.
Thus $\int_{\mathcal D_1}P=\int_{\mathcal D_2}P$.

The two finite-dimensional volumes are
\begin{align}
 \operatorname{vol}(\mathcal F_{\mathrm g})
    &=r_q^{V-1}(\det{}'\Delta_0/V)^{1/2},
                                     \label{eq:external-gauge-volume}\\
 \dd A_{\mathrm h}
    &=V^{1/2}q^{-3}\,\dd\theta
      =r_q^3V^{1/2}\frac{\dd\theta}{(2\pi)^3}.
                                     \label{eq:external-twist-Jacobian}
\end{align}
The first follows from the square root of the Gram
determinant. For the second, the constant cochains
$V^{-1/2}\mathbf1_\mu$ are orthonormal, and the
matrix of $\theta\mapsto H_\theta$ in this basis
has diagonal entries $V^{1/2}/(qn_\mu)$.
Their product is $V^{1/2}q^{-3}$.
Their combined factor is precisely $\mathcal N_n$.
Since $P(A_{\mathrm g}+A_{\mathrm h}+A_{\mathrm t})
=\Phi(dA_{\mathrm t})\mathrm e^{-\|dA_{\mathrm t}\|^2/2}
W(A_{\mathrm t}+A_{\mathrm h})$, it is independent
of $A_{\mathrm g}$. Fubini over the orthogonal
decomposition and the two displayed volume formulas
prove \eqref{eq:external-zero-unfolding}, including
absolute convergence of its right side.

Apply this identity to
$W(A)=\det_n(D_A+\mu_{\mathrm{lat}})$. Its hypotheses
follow from \eqref{eq:external-determinant-invariances}
and \eqref{eq:external-determinant-strip}. If
$Z_n^{(0)}[\Phi]$ is the zero-flux part of the
configuration integral, we have proved
\begin{equation}\label{eq:external-zero-sector}
 Z_n^{(0)}[\Phi]=\mathcal N_n
       \int\frac{\dd\theta}{(2\pi)^3}
       \int_{V_{\mathrm t}}\Phi(dA)\mathrm e^{-\|dA\|^2/2}
          \det_n(D_{A+H_\theta}+\mu_{\mathrm{lat}})\,\dd\sigma(A).
\end{equation}
The twist integral always runs over $[-\pi,\pi]^3$.
For real $A$ and $\theta$ the determinant is
nonnegative. Also $Z_n^{(0)}[1]>0$ at every real
$\mu_{\mathrm{lat}}$: the lower bound in
\cite[Proposition~2.3]{ChamberLattice} was obtained
from the single term $z=0$, which belongs to this
sector. Dividing \eqref{eq:external-zero-sector}
by its value at $\Phi=1$ identifies the zero-flux
curvature law with the law of $dA$ under the
harmonic-averaged transverse measure.

\subsection{The Gaussian normalization and matter density}
\label{subsec:external-Gaussian-law}

Let $F_n=dC^1(n)=dV_{\mathrm t}$ and let $\gamma_n$
be the standard Gaussian probability measure on this
Euclidean space. The map $d_{\mathrm t}=d|_{V_{\mathrm t}}$
is a linear bijection onto $F_n$, by
\eqref{eq:external-Hodge}; write $A(u)=d_{\mathrm t}^{-1}u$.
Changing variables along this map gives
\begin{equation}\label{eq:external-Gaussian-change}
 \int_{V_{\mathrm t}}G(dA)\mathrm e^{-\|dA\|^2/2}\,\dd\sigma(A)
                     =\mathcal G_n\int G(u)\,\dd\gamma_n(u).
\end{equation}
Indeed the Jacobian is $|\det d_{\mathrm t}|^{-1}$,
and $\dim F_n=2V-2$. The Gaussian density contributes
$(2\pi)^{V-1}$; setting $G=1$ identifies the constant
with \eqref{eq:external-transverse-Gaussian}.
The formula holds for nonnegative Borel functions,
and for the exponentially bounded complex functions
used above by their Gaussian integrability.

Define the photon factor and normalized matter factor by
\begin{align}
 Z_n^{\mathrm{ph}}&=\mathcal N_n\mathcal G_n
       =r_q^{V+2}(2\pi)^{V-1}
                            \mathrm e^{-\mathcal L(n)/2},
                                  \label{eq:external-photon-factor}\\
 z_{r,n}(q,\mu_{\mathrm{lat}})
   &=\int\frac{\dd\theta}{(2\pi)^3}
       \int\det_n(D_{A(u)+H_\theta}+\mu_{\mathrm{lat}})
                                             \,\dd\gamma_n(u).
                                  \label{eq:external-matter-factor}
\end{align}
Then \eqref{eq:external-zero-sector} and
\eqref{eq:external-Gaussian-change} imply
\begin{equation}\label{eq:external-zero-factorization}
 Z_n^{(0)}[1]=Z_n^{\mathrm{ph}}z_{r,n},\qquad
                         0<z_{r,n}\le K_n(0).
\end{equation}
Positivity follows from that of $Z_n^{(0)}[1]$;
the upper bound follows because both integration
measures in \eqref{eq:external-matter-factor} are
probability measures.

On a cubic torus $V=n^3$, fix a positive reference mass
$\bar m$, mesh $\eta$ and a real pair
$(m_{\mathrm{in}},\varepsilon_{\mathrm{in}}^{\mathrm v})$,
with $\mu_{\mathrm{lat}}=\bar m\eta+m_{\mathrm{in}}$.
The density with external cochain $a$ is
\begin{equation}\label{eq:external-matter-density}
 \mathrm{dens}^{\theta,a}(u)=
 \mathrm e^{-\varepsilon_{\mathrm{in}}^{\mathrm v}n^3}
 \frac{\det_n(D_{A(u)+H_\theta+a}+\mu_{\mathrm{lat}})}
                       {\det_n(D^{\mathrm{ch}}[0]+\bar m\eta)}.
\end{equation}
The denominator is positive, since the real part
of $D^{\mathrm{ch}}[0]+\bar m\eta$ is at least $\bar m\eta>0$.
At $a=0$ this is a nonnegative bounded continuous
density. Let $\nu^{\theta,0}$ be its measure against
$\gamma_n$ and let $\bar\nu$ be the twist average.
Equations \eqref{eq:external-zero-sector} and
\eqref{eq:external-Gaussian-change} give
\begin{align}
 \bar\nu(\Phi)
   &=\frac{\mathrm e^{-\varepsilon_{\mathrm{in}}^{\mathrm v}n^3}}
        {\det_n(D^{\mathrm{ch}}[0]+\bar m\eta)Z_n^{\mathrm{ph}}}
                                      Z_n^{(0)}[\Phi],
                                  \label{eq:external-density-zero-sector}\\
 \bar\nu(1)
   &=\frac{\mathrm e^{-\varepsilon_{\mathrm{in}}^{\mathrm v}n^3}
                                   z_{r,n}}
                       {\det_n(D^{\mathrm{ch}}[0]+\bar m\eta)}
                                  \in(0,\infty).
                                  \label{eq:external-density-normalization}
\end{align}
Thus $\bar\nu/\bar\nu(1)$ is exactly the zero-flux
curvature law. The vacuum normalization cancels in
this quotient. The quantities $Z_n^{\mathrm{ph}}$
and $z_{r,n}$ depend only on $(n,q,\mu_{\mathrm{lat}})$
and the fixed chamber parameters; the auxiliary
split $\mu_{\mathrm{lat}}=\bar m\eta+m_{\mathrm{in}}$
does not change them.

\subsection{Holomorphic external partition functions}
\label{subsec:external-holomorphy}

Write $\mathfrak F(n)$ for the exponentially bounded
Borel functions $\Phi$ used above. For $w\in\Z^3$ let
$\mathcal M_w$ be the closed integer two-cochains of
flux $w$. For a complex one-cochain $a$ define
\begin{align}
 Z_n^{\mathrm{x},(w)}(a;\Phi)
  &=\sum_{z\in\mathcal M_w}\int_{[0,r_q)^{\mathcal L_n}}
       \Phi(F)\mathrm e^{-\|F\|^2/2}
                \det_n(D_{A+a}+\mu_{\mathrm{lat}})\,\dd A,
       \qquad F=dA+r_qz,\label{eq:external-sector-source}\\
 \mathbf Z_n^{\mathrm{x}}(a;\Phi)
  &=\int\frac{\dd\theta}{(2\pi)^3}
       \int_{V_{\mathrm t}}\Phi(dA)\mathrm e^{-\|dA\|^2/2}
       \det_n(D_{A+H_\theta+a}+\mu_{\mathrm{lat}})\,\dd\sigma(A).
                                  \label{eq:external-averaged-source}
\end{align}
When $\Phi=1$ we omit it from the notation. These are
partition functions in an external connection; the
curvature $F$ in their photon weight is not shifted.

\begin{proposition}\label{prop:external-entire}
For every $\Phi\in\mathfrak F(n)$ both expressions
converge absolutely, locally uniformly in $a$, and are
entire in its $3V$ complex link values. The sum over
all sectors converges absolutely, locally uniformly
as well. Moreover
\begin{equation}\label{eq:external-source-zero-sector}
 Z_n^{\mathrm{x},(0)}(a;\Phi)
             =\mathcal N_n\mathbf Z_n^{\mathrm{x}}(a;\Phi),
 \qquad
 Z_n^{\mathrm{x},(w)}(0;\Phi)=Z_n^{(w)}[\Phi].
\end{equation}
In particular $\mathbf Z_n^{\mathrm{x}}(0)>0$.
\end{proposition}

\begin{proof}
For real $A,\theta$ the imaginary parts of $A+a$ and
$A+H_\theta+a$ are both $\operatorname{Im}a$. On a
polydisc $\max_b|a(b)|\le s$ the determinant is
bounded by $K_n(s)$. If
$|\Phi(u)|\le C_\Phi \mathrm e^{C_\Phi\|u\|}$, the absolute
value of each integrand is consequently bounded by
\[
 K_n(s)C_\Phi \mathrm e^{C_\Phi^2}\mathrm e^{-\|F\|^2/4}
 \quad\text{or}\quad
 K_n(s)C_\Phi \mathrm e^{C_\Phi^2}\mathrm e^{-\|dA\|^2/4}.
\]
The first majorant can be summed over every integer
two-cochain and integrated over the compact box,
giving
$K_n(s)C_\Phi \mathrm e^{C_\Phi^2}(r_qC_q(1/4))^{3V}$.
The second is integrable by
\eqref{eq:external-transverse-Gaussian}: rescaling
$A$ by $\sqrt2$ gives
\[
 \int_{V_{\mathrm t}}\mathrm e^{-\|dA\|^2/4}\,\dd\sigma(A)
                          =2^{V-1}\mathcal G_n.
\]
These bounds also show local uniform absolute
convergence of the sector sum.

We recall the parameter-integral argument to justify
entireness and all later differentiations. Let
$(Y,\nu)$ be a $\sigma$-finite measure space. Suppose
$v(y,a)$ is measurable in $(y,a)$, entire in
$a\in\C^d$ for each $y$, and bounded on each
polydisc by a $\nu$-integrable $M_s(y)$. If $v_\alpha(y)$
are its Taylor coefficients at zero, integration
over the product of circles of radius $s$ gives
\[
 v_\alpha(y)=s^{-|\alpha|}
       \int_{[0,2\pi]^d}
       v(y,s\mathrm e^{i\theta})\mathrm e^{-i\alpha\cdot\theta}
                         \frac{\dd\theta}{(2\pi)^d},
 \qquad |v_\alpha(y)|\le s^{-|\alpha|}M_s(y).
\]
The integral formula makes the coefficients measurable.
For $\max_j|a_j|\le s$, applying the bound at radius
$2s$ yields
\[
 \sum_\alpha\int|v_\alpha(y)a^\alpha|\,\dd\nu(y)
       \le\sum_\alpha2^{-|\alpha|}\int M_{2s}(y)\,\dd\nu(y)
       =2^d\int M_{2s}(y)\,\dd\nu(y)<\infty.
\]
Thus the Taylor series may be integrated term by term,
and its integral is an entire function with coefficients
$\int v_\alpha$. This also proves differentiation under
the integral at zero. Translating the parameter to an
arbitrary $a^0$ proves it there, using the majorant
$M_{s+\|a^0\|_\infty}$. This argument applies to a
counting measure times Lebesgue measure as well as to
the transverse integral.

Apply this argument to the two displayed majorants.
The dependence on $a$ in each integrand is a finite
sum of exponentials of linear forms, by the determinant
polynomial. The other factors are Borel and independent
of $a$, so the hypotheses hold also when $\Phi$ is
merely Borel. This proves the asserted holomorphy.
For fixed complex $a$, the function
$W_a(A)=\det_n(D_{A+a}+\mu_{\mathrm{lat}})$ on real
cochains is bounded by $K_n(\|\operatorname{Im}a\|_\infty)$,
continuous, and has the two invariances required in
\eqref{eq:external-zero-unfolding}. That formula gives
\eqref{eq:external-source-zero-sector}; the value at zero
follows from the definitions. Its positive zero value
follows from $Z_n^{(0)}[1]>0$ and $\mathcal N_n>0$.
\end{proof}

The complex densities in
\eqref{eq:external-matter-density} define finite complex
measures $\nu^{\theta,a}$ and their twist average
$\bar\nu^a$ against $\gamma_n$. Their absolute values
are bounded by
\begin{equation}\label{eq:external-complex-density-bound}
 |\mathrm{dens}^{\theta,a}(u)|
 \le
 \frac{\mathrm e^{-\varepsilon_{\mathrm{in}}^{\mathrm v}n^3}
                   K_n(\|\operatorname{Im}a\|_\infty)}
                   {\det_n(D^{\mathrm{ch}}[0]+\bar m\eta)}.
\end{equation}
This bound is independent of $(\theta,u)$.
Fubini therefore writes the averaged measure as a
bounded Borel density against $\gamma_n$. Applying
\eqref{eq:external-Gaussian-change} at each twist,
and then \eqref{eq:external-source-zero-sector}, gives
\begin{equation}\label{eq:external-complex-measure}
 \bar\nu^a(\Phi)=
 \frac{\mathrm e^{-\varepsilon_{\mathrm{in}}^{\mathrm v}n^3}}
      {\det_n(D^{\mathrm{ch}}[0]+\bar m\eta)\mathcal G_n}
                     \mathbf Z_n^{\mathrm{x}}(a;\Phi)
 =
 \frac{\mathrm e^{-\varepsilon_{\mathrm{in}}^{\mathrm v}n^3}}
      {\det_n(D^{\mathrm{ch}}[0]+\bar m\eta)Z_n^{\mathrm{ph}}}
                     Z_n^{\mathrm{x},(0)}(a;\Phi).
\end{equation}
The exponential majorants in the proof above justify
both changes of integration for $\Phi\in\mathfrak F(n)$.
In particular, for every Borel set $E\subset F_n$,
$a\mapsto\bar\nu^a(E)$ is entire and
\[
 |\bar\nu^a(E)|\le
 \frac{\mathrm e^{-\varepsilon_{\mathrm{in}}^{\mathrm v}n^3}
                   K_n(\|\operatorname{Im}a\|_\infty)}
                   {\det_n(D^{\mathrm{ch}}[0]+\bar m\eta)}.
\]
Indeed $F_n$ is a closed linear subspace, so its
Borel set indicator extends to a bounded Borel
function on $C^2(n)$. Composition with any linear
parameter map $t\mapsto a(t)$ remains entire.
These statements concern the unnormalized complex
measures; the positive normalization at $a=0$ is
\eqref{eq:external-density-normalization}.

\subsection{The source shift and Wick derivatives}
\label{subsec:external-source-shift}

\begin{proposition}\label{prop:external-shift}
For every flux $w$ and complex one-cochain $a$,
\begin{equation}\label{eq:external-shift-sector}
 Z_n^{(w)}[\mathrm e^{\langle da,F\rangle}]
      =\mathrm e^{\langle da,da\rangle/2}Z_n^{\mathrm{x},(w)}(a).
\end{equation}
Both sides, and their sums over $w$, converge
absolutely. Consequently
\begin{equation}\label{eq:external-shift-all-flux}
 Z_n[\mathrm e^{\langle da,F\rangle}]
      =\mathrm e^{\langle da,da\rangle/2}
                    \sum_w Z_n^{\mathrm{x},(w)}(a).
\end{equation}
The quadratic form in the prefactor is bilinear;
there is no complex conjugation of $da$.
\end{proposition}

\begin{proof}
First let $a$ be real. On real cochains define
\[
 P_a(A)=\sum_{z\in\mathcal M_w}
  \mathrm e^{-\|dA+r_qz\|^2/2+\langle da,dA+r_qz\rangle}
                            \det_n(D_A+\mu_{\mathrm{lat}}).
\]
This is a nonnegative Borel function. Adding
$r_q\mathbf1_b$ to $A$ adds $d\mathbf1_b$ to the
integer two-cochain in its curvature. This map
permutes $\mathcal M_w$: its image is closed
and has the same flux, since each exact plaquette
cochain has zero plane sums. The determinant
is unchanged. Hence $P_a$ is invariant under
$r_qC^1(n;\Z)$. Both the compact box and its
translate by $a$ are fundamental domains for
this group. The domain argument in
\cref{subsec:external-zero-flux} gives
\[
 \int_{[0,r_q)^{\mathcal L_n}}P_a(A)\,\dd A
        =\int_{[0,r_q)^{\mathcal L_n}}P_a(A+a)\,\dd A.
\]
On the right the exponent completes to
\[
 -\tfrac12\|F+da\|^2+\langle da,F+da\rangle
                  =-\tfrac12\|F\|^2+\tfrac12\|da\|^2.
\]
Tonelli now identifies the two sides with
\eqref{eq:external-shift-sector} for real $a$.
Their finiteness also follows from the Gaussian
majorant used below.

For complex $a$, the right side of
\eqref{eq:external-shift-sector} is entire by
\Cref{prop:external-entire}.
The left side is entire by the same
parameter-integral argument. In fact, on
$\max_b|a(b)|\le s$,
\[
 \|d\operatorname{Re}a\|\le4s\sqrt{3V}=:c_s,
 \qquad
 \left|\mathrm e^{-\|F\|^2/2+\langle da,F\rangle}
                            \det_n(D_A+\mu_{\mathrm{lat}})\right|
                  \le K_n(0)\mathrm e^{c_s^2}\mathrm e^{-\|F\|^2/4}.
\]
This can be summed and integrated over all integer
two-cochains, so it supplies an integrable polydisc
majorant. Both sides agree on the real cochain
space. An entire function vanishing there has
all complex partial derivatives zero there:
at a real point these derivatives can be taken
along real increments. Repeating this for all
derivatives shows that every Taylor coefficient
at zero vanishes. Thus the two entire functions
agree for all complex $a$.
Finally the same Gaussian majorant permits
summation over sectors; the sector sum on the
right converges absolutely by
\Cref{prop:external-entire}.
\end{proof}

Let $a^{(1)},\ldots,a^{(r)}$ be real cochains and
$X_j=\langle da^{(j)},F\rangle$. Dividing
\eqref{eq:external-shift-sector} at $w=0$ by its
positive zero-source value yields the normalized
generating identity
\begin{equation}\label{eq:external-generating-general}
 \E_n^{\mathrm{av}}\exp\left(
       \sum_jz_jX_j-\tfrac12\sum_{i,j}
                  z_iz_j\langle da^{(i)},da^{(j)}\rangle\right)
 =
 \frac{\mathbf Z_n^{\mathrm{x}}(\sum_jz_ja^{(j)})}
                         {\mathbf Z_n^{\mathrm{x}}(0)},
                                                   \qquad z\in\C^r.
\end{equation}
The expectation $\E_n^{\mathrm{av}}$ is the zero-flux
law identified in \eqref{eq:external-zero-sector}.
The Gaussian factor on the left is determined by
the counting pairing of the curvatures of the
source cochains.

Use the reference cells of \Cref{sec:reference-cells}
on the torus $n_*$. Their index set is
$\mathcal C_*=\{0,\ldots,2L-1\}^3$, and
$\widehat a_c$ is the periodization of
$\hat{\mathfrak a}_c^o$. These are real periodic cochains.
Equation \eqref{eq:cells-orthogonality} gives
\begin{equation}\label{eq:external-cell-pairing}
 \langle d\widehat a_c,d\widehat a_{c'}\rangle
                         =v_\eta\,1_{\{c=c'\}}.
\end{equation}
The variance $v_\eta=\|d\hat{\mathfrak a}_\eta^o\|^2$
is independent of the cell and the placement.
This is the cell-geometry input needed here:
the external-response proof uses neither the
cell width nor its offsets otherwise. Put
\[
 a(t)=e^{-1}\sum_{c\in\mathcal C_*}t_c\widehat a_c,\qquad
 \bar Z^{\mathrm{ext}}(t)=
     \int\frac{\dd\theta}{(2\pi)^3}
        \int\det_{n_*}(D_{A(u)+H_\theta+a(t)}
                              +\mu_{\mathrm{lat}})\,\dd\gamma_{n_*}(u).
\]
The reference torus has $n_*\ge4$, and $e>0$ is
the physical charge in this convention.
Equation \eqref{eq:external-Gaussian-change} gives
\begin{equation}\label{eq:external-reference-partition}
 \bar Z^{\mathrm{ext}}(t)=
          \mathcal G_{n_*}^{-1}\mathbf Z_{n_*}^{\mathrm{x}}(a(t)),
 \qquad
 \mathring\Phi(t):=\frac{\bar Z^{\mathrm{ext}}(t)}
                           {\bar Z^{\mathrm{ext}}(0)}
        =\frac{\mathbf Z_{n_*}^{\mathrm{x}}(a(t))}
                           {\mathbf Z_{n_*}^{\mathrm{x}}(0)}.
\end{equation}
Its zero value in the denominator is positive,
and $\mathring\Phi$ is entire. With
$X_c=\langle d\widehat a_c,F\rangle$,
\eqref{eq:external-generating-general} becomes
\begin{equation}\label{eq:external-generating-cells}
 \E_{n_*}^{\mathrm{av}}
      \exp\sum_c(z_cX_c-\tfrac12v_\eta z_c^2)
                                      =\mathring\Phi(ez).
\end{equation}

Define the Hermite polynomials with this variance by
$H_j(y)=\partial_z^j \mathrm e^{zy-v_\eta z^2/2}|_{z=0}$.
In particular
\[
 H_2(y)=y^2-v_\eta,\qquad
 H_4(y)=y^4-6v_\eta y^2+3v_\eta^2.
\]
To differentiate \eqref{eq:external-generating-cells}
under the expectation, on a polydisc
$\max_c|z_c|\le s$ use
\[
 \left|\exp\sum_c(z_cX_c-\tfrac12v_\eta z_c^2)\right|
 \le
 \exp\left(s|\mathcal C_*|\sqrt{v_\eta}\|F\|
                       +\tfrac12|\mathcal C_*|s^2v_\eta\right).
\]
This is exponentially bounded in the real curvature
and is integrable under the zero-flux law, by the
Gaussian majorant. The parameter-integral argument
therefore applies to this probability measure.
The resulting products of $H_2$ or $H_4$ are also
integrable: $|X_c|\le\sqrt{v_\eta}\|F\|$ bounds them by a
finite polynomial in $\|F\|$, which is bounded by
a constant times $\mathrm e^{\|F\|}$.

For $m=2$ or $4$ and a subset $S\subset\mathcal C_*$,
differentiate in precisely those variables. The
factor $ez$ on the right supplies $e^{m|S|}$, giving
\[
 \E_{n_*}^{\mathrm{av}}\prod_{c\in S}e^{-m}H_m(X_c)
         =\left(\prod_{c\in S}\partial_{t_c}^m
                                      \mathring\Phi\right)(0).
\]
Expanding the product over subsets proves, for
every real $c_m$,
\begin{equation}\label{eq:external-Wick-product}
 \E_{n_*}^{\mathrm{av}}
        \prod_{c\in\mathcal C_*}(e^{-m}H_m(X_c)-c_m)
     =\left(\prod_{c\in\mathcal C_*}
                (\partial_{t_c}^m-c_m)\mathring\Phi\right)(0).
\end{equation}
The derivative operators commute because
$\mathring\Phi$ is entire. This identity is the
finite-volume quantity to which the terminal
interacting estimates will be applied; all its
normalization factors have already been included
in \eqref{eq:external-reference-partition}.


